\section*{第 \ref{ch:g10:relative-motion} 章 --- 相对运动}

\begin{solution}{exo:g10:relative-motion:1}
相对于车厢：静止。相对于地面和树：
以 \qty{100}{km/h} 运动。相对于迎面而来的列车：
相向运动的速率相加，$100 + 100 = \qty{200}{km/h}$。在乘客所乘列车的
参考系中，树沿直线以 \qty{100}{km/h} 的速率
向后运动。
\end{solution}

\begin{solution}{exo:g10:relative-motion:2}
$72 / 3.6 = \qty{20}{m/s}$; $108 / 3.6 = \qty{30}{m/s}$;
$25 \times 3.6 = \qty{90}{km/h}$; $340 \times 3.6 = \qty{1224}{km/h}$.
每组中，以 \unit{km/h} 为单位的数值较大，
这符合数量级检验的要求。
\end{solution}

\begin{solution}{exo:g10:relative-motion:3}
骑车人：$v = \frac{\qty{36}{km}}{\qty{1.5}{h}} = \qty{24}{km/h}
= 24/3.6 \approx \qty{6.7}{m/s}$。短跑运动员：
$v = \frac{\qty{100}{m}}{\qty{10.0}{s}} = \qty{10.0}{m/s}
= \qty{36}{km/h}$。短跑运动员的速率是骑车人的一倍半，但只能持续
十秒，而骑车人能保持这一速率一个半小时。
只有在时长相当时，平均速率的比较才公平。
\end{solution}

\begin{solution}{exo:g10:relative-motion:4}
同向：$1.4 + 0.9 = \qty{2.3}{m/s}$。逆着传送带方向：
$1.4 - 0.9 = \qty{0.5}{m/s}$。站在传送带上静止不动：传送带自身的
速率为 \qty{0.9}{m/s}。
\end{solution}

\begin{solution}{exo:g10:relative-motion:5}
下落的钢笔：用地面参考系也可以，但以\emph{列车}为参考系更好，
在其中钢笔只是竖直下落。空间站：
地心参考系（它整体绕地球运动，忽略地球自转）。
火星年：日心参考系（火星近似沿圆形轨道绕
太阳运行）。短跑：地面参考系（起跑线、终点线和秒表
均固定在地面上）。
\end{solution}

\begin{solution}{exo:g10:relative-motion:6}
\emph{1.} $108/3.6 = \qty{30}{m/s}$。朝前方：
$30 + 1.5 = \qty{31.5}{m/s}$；朝后方：
$30 - 1.5 = \qty{28.5}{m/s}$.

\emph{2.} 相对于列车，乘客以
\qty{1.5}{m/s} 的速率走过 \qty{60}{m}：$t = 60/1.5 = \qty{40}{s}$。相对于地面，乘客
走过的距离为 $31.5 \times 40 = \qty{1260}{m}$——即列车带动的 \qty{1200}{m} 加上
乘客步行的 \qty{60}{m}，
与速度合成规律一致。
\end{solution}

\begin{solution}{exo:g10:relative-motion:7}
\emph{1.} 顺流：$5.0 + 2.0 = \qty{7.0}{m/s}$，故
$t_1 = 420/7.0 = \qty{60}{s}$。逆流：$5.0 - 2.0 = \qty{3.0}{m/s}$，
故 $t_2 = 420/3.0 = \qty{140}{s}$。

\emph{2.} 平均速率：$v = \frac{840}{60 + 140}
= \frac{840}{200} = \qty{4.2}{m/s} < \qty{5.0}{m/s}$。水流的
助力与阻力在\emph{速率}上相等，但在\emph{时间}上并不相等：
船逆流而行的时间为 \qty{140}{s}，顺流仅 \qty{60}{s}，
因此逆流段主导了平均速率。
\end{solution}

\begin{solution}{exo:g10:relative-motion:8}
\emph{1.} 相向运动：在 A 的参考系中，B 以
$25 + 20 = \qty{45}{m/s}$ 的速率靠近。当 B 滑过
两者长度之和 $180 + 90 = \qty{270}{m}$ 时，相遇过程结束：
$t = 270/45 = \qty{6.0}{s}$.

\emph{2.} 同向运动：相对速率为 $25 - 20 = \qty{5.0}{m/s}$，
且 $t = 270/5.0 = \qty{54}{s}$——时间是前者的九倍，这就是
超越卡车时感觉漫长的原因。
\end{solution}

\begin{solution}{exo:g10:relative-motion:9}
以骑车人为参考系：是一个圆（轮缘所画的圆，圆心在轮毂上）。
地面参考系：是摆线——一连串拱形，
车轮每转一圈对应一个拱。以气门嘴本身为参考系：气门嘴静止，
只是一个点。地面参考系中的轮毂：是一条水平直线，
离地高度等于一个车轮的半径。
\end{solution}

\begin{solution}{exo:g10:relative-motion:10}
\emph{1.} $v = \frac{\qty{40000}{km}}{\qty{24}{h}}
\approx \qty{1.7e3}{km/h} = 1667/3.6 \approx \qty{463}{m/s}$ ---
比声速还快。

\emph{2.} 我们感觉不到，是因为周围的一切——地面、
空气、建筑物——都具有完全相同的运动，因此相对于我们，没有任何物体在运动
（关于惯性原理，
见\cref{ch:g10:inertia}，其中说明共同的匀速运动为何在内部无法
察觉）。“太阳升起”是地面参考系下的描述，
“地球自转”则是地心参考系下的描述，二者针对
同一观察结果。
\end{solution}

\begin{solution}{exo:g10:relative-motion:11}
设 $L$ 为自动扶梯的长度。扶梯速率为 $L/60$；
步行速率为 $L/90$。人在运动的扶梯上步行时，速率相加：
\[
v = \frac{L}{60} + \frac{L}{90} = L\left(\frac{3 + 2}{180}\right)
= \frac{L}{36},
\]
故乘行时间为 \qty{36}{s}。（不是 $60$ 与 $90$ 的平均值：
速率相加，时间不能相加。）
\end{solution}

\begin{solution}{exo:g10:relative-motion:12}
\emph{1.} 设单程距离为 $d$：
$t = \frac{d}{60} + \frac{d}{90} = d\,\frac{3 + 2}{180} = \frac{d}{36}$,
故 $v = \frac{2d}{d/36} = \qty{72}{km/h}$。

\emph{2.} 若直接对两个速率取平均得 $75$，
相当于假定两段路上所花的\emph{时间}相等；
事实上两段距离相等，故速率较慢的一段用时更长（$\frac{d}{60} > \frac{d}{90}$），
从而把平均值拉向 \qty{60}{km/h}。平均速率是总路程除以
总时间，仅此而已。
\end{solution}

\begin{solution}{exo:g10:relative-motion:13}
\emph{1.} 以第一名骑车人为参考系，第二名骑车人沿
直线靠近，速率为 $15 + 20 = \qty{35}{km/h}$，起始距离为 \qty{35}{km}：
两人恰好在 \qty{1.0}{h} 后相遇。

\emph{2.} 第一名骑车人骑行 $15 \times 1 = \qty{15}{km}$，第二名
骑行 $20 \times 1 = \qty{20}{km}$ --- 合计 \qty{35}{km}，与预期相符。

\emph{3.} 同向骑行时，追近速率为 $20 - 15 = \qty{5}{km/h}$，
故追及所需时间为 $35/5 = \qty{7}{h}$。
\end{solution}

\begin{solution}{exo:g10:relative-motion:14}
\emph{1.} 水流推动游泳者沿河流方向运动，而非横穿
河流：两个运动独立合成，朝对岸的前进
速率恒为 \qty{1.2}{m/s}，与漂移无关。渡河时间为
$t = 60/1.2 = \qty{50}{s}$.

\emph{2.} 向下游漂移：$0.9 \times 50 = \qty{45}{m}$。

\emph{3.} 地面参考系中的速度是直角三角形的斜边，
两直角边分别为 $1.2$ 和 $0.9$：
$v = \sqrt{1.2^2 + 0.9^2} = \sqrt{2.25} = \qty{1.5}{m/s}$（即 $3$--$4$--$5$
三角形按 $0.3$ 缩放）。路径长度：$1.5 \times 50 = \qty{75}{m}$
--- 亦可直接由 $\sqrt{60^2 + 45^2} = \qty{75}{m}$ 得出。
\end{solution}

\begin{solution}{exo:g10:relative-motion:15}
\emph{1.} $v = \frac{2\pi \times \num{1.50e11}}{\num{3.16e7}}
= \frac{\num{9.42e11}}{\num{3.16e7}} \approx \qty{2.98e4}{m/s}
\approx \qty{30}{km/s}$.

\emph{2.} $30 / 0.46 \approx 65$：地球带着我们绕太阳运动的速度
约为其绕地轴自转速度的 65 倍。

\emph{3.} 在地心参考系中，火星的视路径是其自身公转
与地球运动的合成，而当较快的地球追上火星时，
合成运动会在数周内反向运行。在
日心参考系中，两颗行星都只是以稳定速率绕太阳运动，
--- 所谓“逆行环”并非火星的真实运动，而是
参考系选择造成的假象。
\end{solution}

\begin{solution}{pb:g10:relative-motion:1}
\textbf{1.} $t = 90/1.0 = \qty{90}{s}$.

\textbf{2.} 地面速率为 $1.5 + 1.0 = \qty{2.5}{m/s}$：
$t = 90/2.5 = \qty{36}{s}$.

\textbf{3.} $t = 90/1.5 = \qty{60}{s}$.

\textbf{4.} 相对地面的速率为 $1.5 - 1.0 = \qty{0.5}{m/s}$，
方向与传送带相反：$t = 90/0.5 = \qty{180}{s}$ --- 快步行走
整整三分钟。

\textbf{5.} 地面速率为 $1.0 - 1.0 = \qty{0}{m/s}$：孩子拼命奔跑，
却原地不动，在地面参考系中静止，而在
传送带参考系中却在疾跑。这正是跑步机。

\textbf{6.} 传送带参考系 / 地面参考系：站立的乘客 $0$ /
\qty{1.0}{m/s}；随传送带行走的人 $1.5$ / \qty{2.5}{m/s}；在
地面上行走的人 $0.5$（地面在传送带参考系中向后滑动）/
\qty{1.5}{m/s}；恶作剧者 $1.5$ / \qty{0.5}{m/s}；孩子 $1.0$ /
\qty{0}{m/s}。每人双腿感受到的，是脚下之物所在参考系中的
那一列数值 --- 对传送带上者取传送带参考系，对
地面上行走者取地面参考系：用力程度即为相对脚下之物的
行走速率。

\textbf{7.} 渡河方向的前进速率恒为 \qty{2.0}{m/s}，
与水流无关：$t = 120/2.0 = \qty{60}{s}$。

\textbf{8.} 向下游漂移：$1.5 \times 60 = \qty{90}{m}$。

\textbf{9.} 地面速率是直角三角形的斜边，两直角边分别为
$2.0$ 和 $1.5$：$v = \sqrt{2.0^2 + 1.5^2} = \sqrt{6.25}
= \qty{2.5}{m/s}$。轨迹长度：$2.5 \times 60 = \qty{150}{m}$
--- 又是 $3$--$4$--$5$ 这组勾股数，按比例放大为 $90$--$120$--$150$。

\textbf{10.} 水流参考系中：一条直线，\qty{120}{m}，与
河岸垂直。地面参考系中也是直线，但是倾斜的，且
\qty{150}{m} 长。形状相同（都是直线），方向与
长度不同：轨迹属于参考系，而不属于渡船。

\textbf{11.} 渡河的前进速率为 \qty{1.0}{m/s}，因此
$t = 120/1.0 = \qty{120}{s}$。漂移：$1.5 \times 120 = \qty{180}{m}$
--- 为渡船漂移距离 \qty{90}{m} 的两倍。地面速率为斜边，两直角边分别为
$1.0$ 和 $1.5$：$v = \sqrt{1.0^2 + 1.5^2} = \sqrt{3.25}
\approx \qty{1.8}{m/s}$。速率较慢的游泳者渡河所需时间加倍，
因此同一水流有两倍的时间把他们
冲向下游。

\textbf{12.} 地面参考系中：原木以水流速率
\qty{1.5}{m/s} 漂移。水流参考系中：原木静止不动，随周围的水
一起漂浮。原木是否“运动”，并无脱离参考系的唯一答案 --- 这
正是本章的要点。

\textbf{13.} 取地心参考系。轨道半径
$r = \num{6.37e6} + \num{4.0e5} = \qty{6.77e6}{m}$，因此
\[
v = \frac{2\pi r}{T} = \frac{2\pi \times \num{6.77e6}}{\num{5.56e3}}
\approx \qty{7.65e3}{m/s} \approx \qty{7.7}{km/s}.
\]

\textbf{14.} 空间站完成一圈轨道运行期间，地球
在其下方转过约 $23^\circ$，因此地面上的轨迹
每经过一次便向西偏移。在地面参考系中，空间站的路径是一条
复杂的波浪形曲线，从不精确重复；在地心参考系中，
它是一个固定的近似圆。

\textbf{15.} $7.65 / 0.46 \approx 17$：空间站超过自转赤道的速率约为
赤道速率的约十七倍，这正是它能在
几分钟内划过整个天空的原因。

\textbf{16.} 在地心参考系中，$v = \frac{2\pi \times \num{4.22e7}}{\num{8.64e4}}
\approx \qty{3.07e3}{m/s} \approx \qty{3.1}{km/s}$。
在地面参考系中，它是\emph{静止}的，悬停在赤道上
某一固定点的上方 --- 它绕地轴运行一周的时间
恰好等于地面自转一周的时间，因此在地面参考系中两种运动
相互抵消。静止与 \qty{3.1}{km/s}：两者都对，各属一个参考系。

\textbf{17.} 在日心参考系中，月球的运动基本上
沿着地球的轨道 --- 即绕太阳的一个巨大近似圆 ---
叠加上周期为 \qty{27}{day} 的轻微摆动。太阳的
影响占绝对主导，以至于轨迹始终弯向
太阳，从不出现向后的回绕。

\textbf{18.} “太阳升起”是地面参考系中的描述，其中
观察者静止，太阳掠过天空。“地球自转”
是地心参考系中的描述，其中太阳的
方向稳定，观察者随地球绕
地轴向东运动。两者描述同一观测，各自在其参考系中都正确，
\cref{prop:g10:relative-motion:relativity}说明仅凭运动本身的
观测无法作出裁决。

\textbf{19.} 托勒密的地心参考系自然地与天空的日周运动相符，
并能较好地预测行星位置，但
只能靠层层叠加的本轮来重现每颗
行星的游移。哥白尼的日心参考系使每颗行星都成为
稳定的近似圆，并把火星的逆行回环变成了一种
参考系效应 --- 地球追上了内道上较慢的跑者。
现代的结论并非太阳“真正”静止（它也
绕银河系中心运动），而是参考系只是工具：日心参考系
使行星运动的描述简单得多，而
简洁而非真理，才是参考系的用处。

\textbf{20.} 地面参考系中：\qty{0}{km/s} --- 在扶手椅上静止。
地心参考系中：约 \qty{0.46}{km/s}，随地球绕
地轴运动。日心参考系中：约 \qty{30}{km/s}，随太阳绕行。
观看一部 \qty{90}{min} 的电影时：$d = \num{3.0e4} \times 5400
\approx \qty{1.6e8}{m} = \qty{160000}{km}$ --- 几乎到达
月球距离的一半，人却未离开扶手椅。寓意：运动不是物体的固有属性，
而是物体与参考系的共同结果；应选择使问题最简单的参考系，
并始终说明你所选择的参考系。
\end{solution}
